\section{Del Online Robot Learning a la Offline Recipe}

Esta sección desarrolla uno por uno los tres ingredientes y usa seis años de investigación en robótica de Google Brain para explicar por qué el equipo terminó eligiendo multi-task imitation. El cambio decisivo no fue abandonar el aprendizaje en línea, sino convertir primero los datos en una infraestructura independiente, de modo que varios algoritmos puedan consumir una y otra vez el mismo conjunto de experiencia real.

\subsection{Ingredient 1: Scaling-Friendly Architecture}

El primer ingrediente proviene del scaling en lenguaje y visión: network de alta capacidad, self-attention, más compute/data, tokenización e interoperabilidad entre módulos. El valor del Transformer no reside solo en la fórmula de attention, sino en que ofrece una interfaz de secuencia unificada que permite que el historial de imágenes, las instrucciones en lenguaje y las acciones entren en un mismo autoregressive model.

\lecturefigure{slide-07-ingredient-ml-scaling.jpg}{Ingredient 1: tomar de ML scaling la architecture, el compute, los data y la tokenización}{Diapositiva V007 recuperada del video oficial; Stanford CS25 Lecture 15}

\begin{knowledgebox}{Qué significa la tokenización en robótica}
En NLP, un token suele ser una subpalabra discreta; en RT-1, la imagen se convierte primero en tokens visuales compactos y las magnitudes de control continuas se cuantizan en action tokens. Una token interface unificada permite reutilizar el Transformer, pero el error de cuantización, la longitud de la secuencia y el costo de decodificación en tiempo real se vuelven nuevas restricciones de ingeniería.
\end{knowledgebox}

\subsection{Ingredient 2: aprovechar Internet-Scale Models}

El segundo ingrediente no exige que el equipo de robótica entrene desde cero la semántica de lenguaje o de visión. Los modelos externos siguen mejorando en su propio dominio y pueden aportar common sense, comprensión del lenguaje, categorías visuales y representaciones componibles. Para la robótica, estos modelos funcionan mejor como planner, captioner, labeler o representation provider, y no como fuente directa de motor commands cuando carecen de grounding.

\lecturefigure{slide-08-ingredient-internet-models.jpg}{Ingredient 2: aprovechar los Internet-scale generative models, que mejoran continuamente}{Diapositiva V008 recuperada del video oficial; Stanford CS25 Lecture 15}

\teachervoice{El ponente llama a esto «Bitter Lesson 2.0»: no basta con elegir algoritmos que mejoren con más compute; también hay que diseñar la interfaz del robot de modo que una foundation model futura y más potente pueda conectarse directamente. Los modelos de lenguaje actuales no son el punto final, y el sistema debe beneficiarse de los avances de los modelos de mañana.}

\subsection{Ingredient 3: la Robotics pasa de Online a Offline}

El tercer ingrediente es el régimen de datos. En robótica, el RL en línea exige aprender mientras se ejecuta, y los fallos pueden dañar el equipo o desperdiciar tiempo de operación; el offline learning, en cambio, fija primero el conjunto de datos y luego entrena y compara varios modelos con las mismas trajectories. La slide 9 contrapone corpus fuera de línea como The Pile y LAION con el aprendizaje online / on-policy tradicional de la robótica, para subrayar la diferencia en la forma de consumir los datos.

\lecturefigure{slide-09-ingredient-offline-data.jpg}{Ingredient 3: las foundation models consumen corpus fuera de línea gigantescos, mientras que la robótica ha dependido durante mucho tiempo de la recolección en línea}{Diapositiva V009 recuperada del video oficial; Stanford CS25 Lecture 15}

Sean los datos de demostración fuera de línea $\mathcal{D}=\{(o_t,l,a_t)\}$, donde $o_t$ es la observación, $l$ es la tarea expresada en lenguaje y $a_t$ es la acción del experto. El objetivo más directo de behavioral cloning es
\begin{equation}
\theta^*=\arg\min_{\theta}
-\mathbb{E}_{(o_t,l,a_t)\sim\mathcal{D}}
\log p_{\theta}(a_t\mid o_{\le t},l).
\end{equation}
Esto convierte el aprendizaje de control en aprendizaje supervisado, que es estable y fácil de escalar; pero en cuanto la policy se desvía de la trayectoria del experto, encuentra estados que la distribución de entrenamiento no cubre, y esa es la limitación fundamental de la offline imitation.

\subsection{Seis años de historia: de End-to-End a Multi-Task}

La slide histórica conecta el goal conditioning, QT-Opt, sim2real y el RL a gran escala con robots reales de 2016--2020 con los sistemas multitarea de 2020--2022, como BC-Z, IL+RL y MT-Opt. No es una revisión exhaustiva, sino la evolución de las preguntas internas del equipo: primero, «¿puede un robot de extremo a extremo aprender una tarea?»; después, «¿cómo escalar a muchas tareas en escenarios más complejos?».

\lecturefigure{slide-10-google-robotics-history.jpg}{Historia de la robotics en Google Brain: del aprendizaje de extremo a extremo de una sola tarea a los datos y las políticas multitarea}{Diapositiva V010 recuperada del video oficial; Stanford CS25 Lecture 15}

\lecturefigure{slide-11-online-to-offline-shift.jpg}{El giro de seis años: de los online methods a la offline multi-task imitation}{Diapositiva V011 recuperada del video oficial; Stanford CS25 Lecture 15}

\teachervoice{Cuando un estudiante preguntó si una «24/7 robot farm» sería suficiente, el ponente subrayó que las horas de robot no son la única escala. Si todo lo que se recolecta durante el día es el mismo tipo de acción y de escenario, la cantidad de datos crece sin ampliar la task diversity; lo verdaderamente escaso es cubrir nuevos failure modes y variaciones combinatorias.}

\begin{knowledgebox}{El lugar de BC-Z en el curso}
BC-Z usa lenguaje o videos de personas como task conditioning, entrena una multi-task imitation policy en alrededor de cien manipulation tasks y evalúa zero-shot task combinations. Demostró que el templated language basta como interfaz de tareas y también expuso la sensibilidad de la ResNet policy a la distribución de entrenamiento y a escenarios nuevos, con lo que proporcionó a RT-1 su punto de partida en datos y en arquitectura.
\end{knowledgebox}

\subsection{Convertir los Ingredients en una Recipe}

La receta final alinea las tres columnas de ingredientes con las lessons: architecture / tokenización de alta capacidad; external models reemplazables que mejoran continuamente; y datos robóticos fuera de línea abundantes y diversos. La frase inferior los condensa en «large diverse offline datasets + high-capacity architectures + language as universal glue».

\lecturefigure{slide-12-foundation-recipe.jpg}{Turning Ingredients into a Recipe: la combinación de arquitectura, modelos externos y datos fuera de línea}{Diapositiva V012 recuperada del video oficial; Stanford CS25 Lecture 15}

\begin{importantbox}{Desacoplar Data Generation y Data Consumption}
Los sistemas en línea suelen adaptar el proceso de recolección a un algoritmo concreto, y cuando el algoritmo cambia hay que volver a operar los robots; la offline recipe convierte las trayectorias en un activo de largo plazo, de modo que RT-1, el RL futuro o distintas architectures puedan consumir los mismos datos. El desacoplamiento no elimina la distribution shift, pero reduce de forma notable el costo de iteración y mejora la reproducibilidad.
\end{importantbox}

\subsection{Resumen de la sección}

El núcleo de la recipe de la clase no es que «el Transformer lo resuelve todo», sino la interfaz unificada y el régimen de datos: la tokenización permite que varias modalidades entren en un mismo modelo, las foundation models externas aportan un prior reemplazable y los offline datasets desacoplan la recolección del entrenamiento. A continuación, RT-1 llevará estos principios a un presupuesto estricto de control en tiempo real.
